%% ============================================================
%% CHAPTER 7 — Skills Without an Operating System
%%                                        (budget: ~8,000 words)
%% Role: answers RQ3. The design-space chapter.
%% ============================================================
\chapter{Skills Without an Operating System}
\label{ch:skills}

\section{What Linux Skills Silently Assume}
\label{sec:skill-assumptions}
% Dissect OpenClaw/Claude Code skills: executable scripts assume a shell
% and interpreters; isolation assumes processes/containers; dependencies
% assume a package manager; installation assumes a hierarchical FS.
% Remove the OS and each assumption fails — enumerate them as a table
% that the rest of the chapter answers.

\section{A Design Space for OS-less Skills}
\label{sec:skill-space}
% The expressiveness ladder with the OS-substitute column:
% L0 prompt-only (interpreter = the cloud LLM itself);
% L1 declarative config (fixed firmware mechanisms interpret data);
% L2 interpreted script (embedded interpreter = sandbox; timeout hook =
%    process killer; module allowlist = permission system; interpreter
%    heap cap = memory isolation);
% L3 loadable native code (discussed, not built: essentially
%    unsandboxable without an MMU — argues L2 is the sweet spot).
% Axes: expressiveness, footprint, safety, LLM-authorability, lifecycle.

\section{Case Study: ESP-Claw's Lua Skill System}
\label{sec:skill-espclaw}
% Precise description from the public spec: SKILL.md JSON frontmatter,
% cap_groups activation, the single lua_run_script tool, arg_schema
% device-side validation, the async job system (named jobs, exclusive
% locks). Engineering precedent, no published evaluation.

\section{Implementation: MicroPython Skills for MimiClaw}
\label{sec:skill-impl}
% Wiring the vendored micropython_embed component: a python_run_script
% tool mirroring cap_lua; sandbox design (timeout vs the hardware
% watchdog — the bare-metal-specific interaction; module allowlist;
% heap cap); skills as SPIFFS directories; how the model authors a new
% skill at runtime through existing file tools.

\section{Evaluation}
\label{sec:skill-eval}
% Four experiments:
% 1. Coverage: 20 capability-extension requests; % achievable at L0 vs L2.
% 2. LLM authorability: model writes the skill (L0 markdown; L2 script);
%    first-try success and repair iterations, by model tier.
% 3. The interpretation tax: flash/RAM cost of the interpreter; script vs
%    native C tool latency for identical operations.
% 4. Sandbox verification: runaway loop -> timeout kill; allocation bomb
%    -> heap cap; forbidden import -> blocked; WDT interplay.
% Plus the comparative table: Lua (ESP-Claw) vs MicroPython (MimiClaw)
% vs C-only, on footprint/latency/safety mechanisms.

\section{Discussion}
\label{sec:skill-discussion}
% When each level is the right choice; the security surface an
% LLM-authored script opens (link forward to guards, back to Ch3);
% lifecycle gaps (versioning, revocation) as future work.
